\section{A Unified Static+Dynamic Design (Proposal)}
\label{sec:unified}

% LOCKED status: PROPOSAL ONLY in this (v1) manuscript. Present the architecture
% + a security argument under the Sec. 2 threat model. EVALUATION IS FUTURE WORK
% and gates the PoPETs v2 submission (where this section becomes the headline,
% built + measured head-to-head vs static-only and TEE-only).
% Status: NOT YET DESIGNED — this section is a forward-looking sketch.

\dnote{This is the SoK's "so what / our recommendation". Do NOT make measured
claims here in v1 — the design isn't built. Mark evaluation explicitly as future
work to avoid over-claiming.}

\subsection{Motivation from the Evidence}
% Sec. 6 shows static obfuscation breaks and the TEE split holds but costs overhead.
% Motivate a design that keeps static obfuscation's low overhead where it is safe
% and falls back to the dynamic (TEE) path only where attacks bite.
\needswork{argue the design from the holds/breaks evidence}

\subsection{Architecture (Proposed)}
\needswork{proposed unified static+dynamic architecture}

\subsection{Security Argument}
% Reason about each attack class from Sec. 6 under the proposed design + threat model.
\needswork{security argument; NO measured numbers in v1}

\subsection{Evaluation Plan (Future Work)}
% The v2 campaign: build it, measure head-to-head vs static-only (AloePri) and
% TEE-only on the Sec. 6 substrate.
\needswork{state the v2 evaluation plan explicitly as future work}
